% ============================================================
\chapter{Th\'eor\`emes de Convergence}
\label{ch:convergence}
% ============================================================

\section{Introduction}

Voici la question que tout analyste finit par se poser~: quand peut-on \'echanger une limite et une int\'egrale~? L'int\'egrale de Riemann offre une r\'eponse spartiate~: sous convergence uniforme, et c'est tout. Mais dans la pratique --- s\'eries de Fourier, th\'eorie des probabilit\'es, physique statistique --- la convergence uniforme est un luxe rare. C'est pr\'ecis\'ement ici que l'int\'egrale de Lebesgue r\'ev\`ele toute sa puissance. Entre 1902 et 1910, Lebesgue, puis Beppo Levi et Pierre Fatou, \'etablissent trois th\'eor\`emes de convergence qui forment le c\oe ur battant de l'int\'egration moderne. Le th\'eor\`eme de convergence monotone de Beppo Levi (1906) traite les suites croissantes~; le lemme de Fatou (1906) fournit une in\'egalit\'e universelle~; et le th\'eor\`eme de convergence domin\'ee de Lebesgue couronne l'\'edifice en autorisant l'\'echange limite-int\'egrale d\`es qu'une fonction int\'egrable \og domine \fg{} la suite. Ces trois r\'esultats sont, sans exag\'eration, les outils les plus utilis\'es de toute l'analyse moderne.

\section{Th\'eor\`eme de convergence monotone (Beppo Levi)}

\begin{theoreme}[Convergence monotone]
\label{thm:convergence_monotone}
Soit $(f_n)_{n \in \N}$ une suite \emph{croissante} de fonctions mesurables
positives~: $0 \leq f_1 \leq f_2 \leq \ldots$ Soit $f = \lim_n f_n =
\sup_n f_n$. Alors~:
\[
\int_X f\, d\mu = \lim_{n \to \infty} \int_X f_n\, d\mu
= \sup_{n \in \N} \int_X f_n\, d\mu.
\]
\end{theoreme}

\begin{proof}
Posons $I = \lim_n \int f_n\, d\mu$ (la limite existe dans $[0, +\infty]$ car la
suite $(\int f_n\, d\mu)$ est croissante par monotonie de l'int\'egrale).

\emph{In\'egalit\'e $\int f\, d\mu \geq I$}~: comme $f_n \leq f$ pour tout $n$,
$\int f_n\, d\mu \leq \int f\, d\mu$, donc $I \leq \int f\, d\mu$.

\emph{In\'egalit\'e $\int f\, d\mu \leq I$}~: soit $\varphi$ une fonction simple
mesurable avec $0 \leq \varphi \leq f$ et $\alpha \in (0, 1)$. Posons
$E_n = \{x : f_n(x) \geq \alpha \varphi(x)\}$. Alors $E_n \uparrow X$ (car
pour tout $x$, $f_n(x) \uparrow f(x) \geq \varphi(x) > \alpha \varphi(x)$
\'eventuellement). On a~:
\[
\int f_n\, d\mu \geq \int_{E_n} f_n\, d\mu \geq \alpha \int_{E_n} \varphi\, d\mu.
\]
En faisant $n \to \infty$, par continuit\'e monotone de la mesure
$A \mapsto \int_A \varphi\, d\mu$~:
\[
I \geq \alpha \int_X \varphi\, d\mu.
\]
En faisant $\alpha \to 1^-$~: $I \geq \int_X \varphi\, d\mu$. Prenant le supremum
sur toutes les fonctions simples $\varphi \leq f$~: $I \geq \int f\, d\mu$.
\end{proof}

\begin{formules}[title=\textbf{Convergence monotone}]
Si $0 \leq f_n \uparrow f$, alors
$\displaystyle\int f_n\, d\mu \uparrow \int f\, d\mu$.
\end{formules}

\begin{corollaire}[S\'eries de fonctions positives]
Soit $(g_n)_{n \in \N}$ une suite de fonctions mesurables positives. Alors~:
\[
\int_X \sum_{n=0}^{\infty} g_n\, d\mu = \sum_{n=0}^{\infty} \int_X g_n\, d\mu.
\]
\end{corollaire}

\begin{proof}
Appliquer le th\'eor\`eme de convergence monotone \`a $f_N = \sum_{n=0}^N g_n$.
\end{proof}

\begin{corollaire}[Additivit\'e de l'int\'egrale]
Pour $f, g : X \to [0, +\infty]$ mesurables~:
$\int (f + g)\, d\mu = \int f\, d\mu + \int g\, d\mu$.
\end{corollaire}

\begin{proof}
Soit $(\varphi_n)$ et $(\psi_n)$ des suites croissantes de fonctions simples
positives avec $\varphi_n \uparrow f$ et $\psi_n \uparrow g$. Alors
$\varphi_n + \psi_n \uparrow f + g$ et
$\int (\varphi_n + \psi_n)\, d\mu = \int \varphi_n\, d\mu + \int \psi_n\, d\mu$.
Par le th\'eor\`eme de convergence monotone appliqu\'e aux deux membres.
\end{proof}

\section{Lemme de Fatou}

\begin{theoreme}[Lemme de Fatou]
\label{thm:fatou}
Soit $(f_n)_{n \in \N}$ une suite de fonctions mesurables positives. Alors~:
\[
\int_X \liminf_{n \to \infty} f_n\, d\mu \leq
\liminf_{n \to \infty} \int_X f_n\, d\mu.
\]
\end{theoreme}

\begin{proof}
Posons $g_n = \inf_{k \geq n} f_k$. Alors $g_n$ est mesurable positive,
$g_n \leq f_n$ et $g_n \uparrow \liminf_n f_n$. Par le th\'eor\`eme de convergence
monotone~:
\[
\int \liminf_n f_n\, d\mu = \lim_n \int g_n\, d\mu \leq
\liminf_n \int f_n\, d\mu,
\]
la derni\`ere in\'egalit\'e venant de $\int g_n\, d\mu \leq \int f_n\, d\mu$.
\end{proof}

\begin{attention}[title=\textbf{L'in\'egalit\'e peut \^etre stricte}]
Prendre $f_n = n \cdot \mathbf{1}_{(0, 1/n)}$ sur $(\R, \lambda)$. Alors
$f_n \to 0$ p.p. mais $\int f_n\, d\lambda = 1$ pour tout $n$. Donc
$0 = \int \liminf f_n\, d\lambda < \liminf \int f_n\, d\lambda = 1$.
\end{attention}

\begin{corollaire}[Lemme de Fatou inverse]
Si $(f_n)$ est une suite de fonctions mesurables positives et s'il existe
$g \in L^1(\mu)$ telle que $f_n \leq g$ p.p. pour tout $n$, alors~:
\[
\limsup_{n \to \infty} \int_X f_n\, d\mu \leq
\int_X \limsup_{n \to \infty} f_n\, d\mu.
\]
\end{corollaire}

\begin{proof}
Appliquer le lemme de Fatou \`a la suite $g - f_n \geq 0$.
\end{proof}

\section{Th\'eor\`eme de convergence domin\'ee (Lebesgue)}

\begin{theoreme}[Convergence domin\'ee]
\label{thm:convergence_dominee}
Soit $(f_n)$ une suite de fonctions mesurables $X \to \overline{\R}$ telle que~:
\begin{enumerate}
    \item[(i)] $f_n \to f$ $\mu$-p.p. (convergence ponctuelle presque partout)\,;
    \item[(ii)] il existe $g \in L^1(\mu)$, $g \geq 0$, telle que
    $\abs{f_n} \leq g$ $\mu$-p.p. pour tout $n$ (\emph{domination}).
\end{enumerate}
Alors $f \in L^1(\mu)$, chaque $f_n \in L^1(\mu)$, et~:
\[
\lim_{n \to \infty} \int_X f_n\, d\mu = \int_X f\, d\mu.
\]
De plus, $\int_X \abs{f_n - f}\, d\mu \to 0$.
\end{theoreme}

\begin{proof}
Comme $\abs{f_n} \leq g$ p.p. et $f_n \to f$ p.p., on a $\abs{f} \leq g$ p.p.,
donc $f \in L^1(\mu)$.

Consid\'erons les fonctions $g + f_n \geq 0$ et $g - f_n \geq 0$. Par le lemme
de Fatou~:
\begin{align*}
\int (g + f)\, d\mu &\leq \liminf_n \int (g + f_n)\, d\mu
= \int g\, d\mu + \liminf_n \int f_n\, d\mu, \\
\int (g - f)\, d\mu &\leq \liminf_n \int (g - f_n)\, d\mu
= \int g\, d\mu - \limsup_n \int f_n\, d\mu.
\end{align*}
Comme $\int g\, d\mu < \infty$, on en d\'eduit~:
\[
\int f\, d\mu \leq \liminf_n \int f_n\, d\mu \leq
\limsup_n \int f_n\, d\mu \leq \int f\, d\mu.
\]

Pour la derni\`ere assertion, appliquer la convergence domin\'ee \`a
$\abs{f_n - f} \leq 2g$ et $\abs{f_n - f} \to 0$ p.p.
\end{proof}

\begin{formules}[title=\textbf{Th\'eor\`eme de convergence domin\'ee}]
\[
\boxed{f_n \to f \text{ p.p.}, \quad \abs{f_n} \leq g \in L^1
\implies \int f_n\, d\mu \to \int f\, d\mu.}
\]
\end{formules}

\section{Applications}

\begin{corollaire}[Continuit\'e sous le signe int\'egral]
\label{cor:continuite_integrale}
Soit $f : X \times I \to \R$ o\`u $I \subset \R$ est un intervalle. Supposons~:
\begin{enumerate}
    \item pour tout $t \in I$, $x \mapsto f(x, t)$ est mesurable\,;
    \item pour $\mu$-p.p. $x$, $t \mapsto f(x, t)$ est continue en $t_0$\,;
    \item il existe $g \in L^1(\mu)$ telle que $\abs{f(x, t)} \leq g(x)$ pour
    tout $t \in I$ et $\mu$-p.p. $x$.
\end{enumerate}
Alors $F(t) = \int_X f(x, t)\, d\mu(x)$ est continue en $t_0$.
\end{corollaire}

\begin{proof}
Si $t_n \to t_0$, alors $f(x, t_n) \to f(x, t_0)$ p.p. avec $\abs{f(x, t_n)}
\leq g(x)$. Par convergence domin\'ee, $F(t_n) \to F(t_0)$.
\end{proof}

\begin{corollaire}[D\'erivation sous le signe int\'egral]
\label{cor:derivation_integrale}
Soit $f : X \times I \to \R$ o\`u $I$ est un intervalle ouvert. Supposons~:
\begin{enumerate}
    \item pour tout $t \in I$, $x \mapsto f(x, t)$ est int\'egrable\,;
    \item pour $\mu$-p.p. $x$, $t \mapsto f(x, t)$ est d\'erivable sur $I$\,;
    \item il existe $g \in L^1(\mu)$ telle que
    $\abs{\frac{\partial f}{\partial t}(x, t)} \leq g(x)$ pour tout $t \in I$
    et $\mu$-p.p. $x$.
\end{enumerate}
Alors $F(t) = \int_X f(x, t)\, d\mu(x)$ est d\'erivable sur $I$ et~:
\[
F'(t) = \int_X \frac{\partial f}{\partial t}(x, t)\, d\mu(x).
\]
\end{corollaire}

\begin{proof}
Pour $h \neq 0$, $\frac{F(t+h) - F(t)}{h} = \int \frac{f(x, t+h) - f(x, t)}{h}
\, d\mu(x)$. Par le th\'eor\`eme des accroissements finis,
$\abs{\frac{f(x, t+h) - f(x, t)}{h}} \leq g(x)$. En faisant $h \to 0$, la
convergence domin\'ee donne le r\'esultat.
\end{proof}

\section{Convergence dans $L^1$}

\begin{theoreme}[Caract\'erisation de la convergence $L^1$]
Soit $f_n, f \in L^1(\mu)$. Les conditions suivantes sont \'equivalentes~:
\begin{enumerate}
    \item $\int \abs{f_n - f}\, d\mu \to 0$ (convergence $L^1$).
    \item $f_n \to f$ en mesure et la famille $(f_n)$ est uniform\'ement
    int\'egrable.
\end{enumerate}
\end{theoreme}

\begin{definition}[Int\'egrabilit\'e uniforme]
Une famille $(f_i)_{i \in I}$ de fonctions int\'egrables est \emph{uniform\'ement
int\'egrable} si~:
\[
\lim_{M \to \infty} \sup_{i \in I} \int_{\{\abs{f_i} > M\}} \abs{f_i}\, d\mu = 0.
\]
\end{definition}

\begin{theoreme}[Vitali]
\label{thm:vitali_convergence}
Soit $\mu(X) < \infty$. Si $(f_n)$ est uniform\'ement int\'egrable et $f_n \to f$
en mesure, alors $f \in L^1(\mu)$ et $\int \abs{f_n - f}\, d\mu \to 0$.
\end{theoreme}

\section{Exercices}

\begin{exercice}
Calculer $\lim_{n \to \infty} \int_0^{\infty}
\frac{n \sin(x/n)}{x(1 + x^2)}\, dx$ en justifiant l'\'echange limite-int\'egrale.
\end{exercice}

\begin{exercice}
Montrer que $F(\alpha) = \int_0^{\infty} \frac{e^{-\alpha x}}{1 + x^2}\, dx$ est
de classe $C^\infty$ sur $(0, \infty)$ et calculer $F'(\alpha)$.
\end{exercice}

\begin{exercice}
Soit $f \in L^1(\R)$. Montrer que $F(t) = \int_{\R} e^{itx} f(x)\, dx$ est
continue sur $\R$ (c'est la transform\'ee de Fourier de $f$).
\end{exercice}

\begin{exercice}
Donner un exemple de suite $(f_n)$ dans $L^1([0,1], \lambda)$ telle que
$f_n \to 0$ p.p. mais $\int f_n\, d\lambda \not\to 0$. Pourquoi le th\'eor\`eme
de convergence domin\'ee ne s'applique-t-il pas~?
\end{exercice}

\begin{exercice}
Calculer $\lim_{n \to \infty} \int_0^1 (1 + x/n)^n e^{-2x}\, dx$.
\end{exercice}

\begin{exercice}
Soit $f \in L^1(\R)$. Montrer que
$\lim_{h \to 0} \int_{\R} \abs{f(x + h) - f(x)}\, dx = 0$
(continuit\'e de la translation dans $L^1$).
\end{exercice}
